% ===== Appendix: corrections to the book, what we added, sources (after Questions; not presented) =====
\appendix
\partaccent{lavender}\gdef\idiottopicemoji{books}\gdef\idiottopicname{Հավելված}
\newcolumntype{L}[1]{>{\raggedright\arraybackslash}p{#1}}

\begin{frame}{Հավելված Ա. Ուղղումներ գրքի նկատմամբ (1/2)}
  \scriptsize
  \setlength{\tabcolsep}{4pt}
  \renewcommand{\arraystretch}{1.2}
  \begin{tabular}{@{}L{2.3cm}L{3.3cm}L{4.7cm}L{3.0cm}@{}}
    \toprule
    \textbf{Թեմա} & \textbf{Գրքում (2016)} & \textbf{Այսօր} & \textbf{Աղբյուր}\\
    \midrule
    Առաջին զգայարանը & Արգանդում առաջինը հոտառությունն է զարգանում & Սովորաբար առաջինը հպումն է (մոտ 8-րդ շաբաթից), քիմիականները՝ ավելի ուշ & Lecanuet \& Schaal, 1996\\
    Գանգուղեղային նյարդեր & «12 դիմային նյարդ» (facial nerves) & 12 գանգուղեղային նյարդ (cranial nerves). դիմայինը դրանցից 7-րդն է & Դասագրքային\\
    Հոտառական գեներ & «Գենոմի 3\%-ը» & Մոտ 1000 գեն՝ \textbf{գեների} մոտ 3\%-ը. 1991-ի աշխատանքը առնետների վրա էր, Նոբել՝ 2004 & Buck \& Axel, 1991\\
    Համային պտկիկներ & «Պտկիկների հինգ տեսակ» & Հինգ տեսակի ընկալիչ բջիջ, ամեն պտկիկում՝ բոլորը & Chandrashekar et al., 2006\\
    Թմբկաթաղանթ & «Կազմված է երեք ոսկրիկից» & Թաղանթ է. ոսկրիկները նրա հետևում են, միջին ականջում & Դասագրքային\\
    \bottomrule
  \end{tabular}
\end{frame}

\begin{frame}{Հավելված Ա. Ուղղումներ գրքի նկատմամբ (2/2)}
  \scriptsize
  \setlength{\tabcolsep}{4pt}
  \renewcommand{\arraystretch}{1.2}
  \begin{tabular}{@{}L{2.3cm}L{3.3cm}L{4.7cm}L{3.0cm}@{}}
    \toprule
    \textbf{Թեմա} & \textbf{Գրքում (2016)} & \textbf{Այսօր} & \textbf{Աղբյուր}\\
    \midrule
    Հյուբել ու Վիզել & Կողմնորոշման սյուները գտել են 1981-ին & Գտել են 1959-1962-ին, 1981-ը Նոբելյան մրցանակի տարին է & Nobel Prize, 1981\\
    Խորության նշան & Հեռացող ճանապարհը՝ «convergence» & Դա գծային հեռանկար է. convergence-ը աչքերի ներս թեքվելն է & Դասագրքային\\
    Գոտկային կեղև & «Քունքային բլթի խորքում» & Կիսագնդերի ներքին մակերեսին, կոշտ մարմնի վրա & Bush, Luu \& Posner, 2000\\
    Overt / covert & Overt՝ նպատակով, covert՝ ինքնաբերաբար & Overt՝ աչքերը շարժելով, covert՝ առանց. նպատակով/ինքնաբերը այլ բաժանում է & Posner, 1980\\
    Օձեր ու սարդեր & Նկարները՝ 1/300 վ-ից կարճ & Դասական փորձերում՝ 30 մվ (մոտ 1/33 վ), հետո ծածկող պատկեր. «1/300»-ը չհաստատեցինք & Öhman \& Soares, 1994\\
    \bottomrule
  \end{tabular}
\end{frame}

\begin{frame}{Հավելված Բ. Ինչ ենք ավելացրել գրքից դուրս}
  \begin{columns}[T]
    \begin{column}{0.48\textwidth}
      \begin{card}[colback=mint!30!cardwhite, fontupper=\scriptsize]
        {\blackfont Թվեր ու հետազոտություններ}\par\smallskip
        \begin{itemize}\setlength{\itemsep}{0.3mm}
          \item Երկու կետի շեմերը՝ Bates' Guide, 2007
          \item «Պարմեզան» ու «փսխուք»՝ նույն փորձից (Herz \& von Clef, 2001)
          \item Ներկած սպիտակ գինին՝ Morrot et al., 2001
          \item Դռան փորձի թվերը՝ Simons \& Levin, 1998
          \item 1000 գեն ու Նոբել 2004. մոր կերածի հոտը՝ Schaal et al., 2000
          \item Մարսի «դեմքը»՝ NASA/JPL, Viking 1 ու HiRISE
        \end{itemize}
      \end{card}
    \end{column}
    \begin{column}{0.48\textwidth}
      \begin{card}[colback=peach!35!cardwhite, fontupper=\scriptsize]
        {\blackfont Մերն է, ոչ թե գրքից}\par\smallskip
        \begin{itemize}\setlength{\itemsep}{0.3mm}
          \item Կոնֆետի, ֆովեայի, աչքերի, բութ մատի, Պոզների ու փոփոխության փորձերը, «Փորձիր տանը»
          \item Ադելսոնի շախմատային իլյուզիան (1995)
          \item Ամանորը (գրքում՝ նոյեմբերի 5 ու հուլիսի 4)
          \item Գծագրերը սխեմատիկ են
          \item Լուսանկարներ՝ Wikimedia Commons (Adelson/Pbroks13, Mpj29, TeWeBs, Pierrecolantidot, S. Viinamäki, Mbejger), NASA
          \item Մեմերը. բովանդակությունը գրքից է
        \end{itemize}
      \end{card}
    \end{column}
  \end{columns}
\end{frame}

\begin{frame}{Հավելված Գ. Աղբյուրներ. գլուխ 5}
  \scriptsize
  Գրքի ծանոթագրություններից ու տեքստից, որոնց վրա հիմնված են սլայդները\par\smallskip
  \begin{columns}[T]
    \begin{column}{0.48\textwidth}
      \begin{itemize}\setlength{\itemsep}{0.2mm}
        \item Bushdid C. et al. 1 trillion olfactory stimuli. Science, 2014
        \item Gerkin R. C., Castro J. B. The number of olfactory stimuli... is still unknown. eLife, 2015
        \item Buck L., Axel R. Odorant receptors. Cell, 1991
        \item Herz R. S., von Clef J. Verbal labeling and odors. Perception, 2001
        \item Hodgson R. T. 13 US wine competitions. J. Wine Econ., 2009
        \item Auvray M., Spence C. The multisensory perception of flavor. Conscious. Cogn., 2008
        \item Frenzel H. et al. Mechanosensory traits in humans. PLOS Biol., 2012
        \item Hubel D. H., Wiesel T. N. Brain mechanisms of vision. Sci. Am., 1979
      \end{itemize}
    \end{column}
    \begin{column}{0.48\textwidth}
      \begin{itemize}\setlength{\itemsep}{0.2mm}
        \item Cherry E. C. Speech with one and with two ears. JASA, 1953
        \item Kahneman D. Attention and Effort. 1973
        \item Strayer D. L. et al. Cognitive distraction in the automobile. AAA, 2013 (գիրքը՝ AAA-ի Traffic Safety Culture Index, 2013)
        \item Mesgarani N., Chang E. F. Attended speaker. Nature, 2012 (գիրքը՝ Mesgarani et al., Science, 2014)
        \item Mallan K. M. et al. Slithering snakes, angry men. Cogn. Emot., 2013
        \item Simons D. J., Levin D. T. Failure to detect changes to people. Psychon. Bull. Rev., 1998
        \item McCann R. S. et al. Heads-up displays. 1993
      \end{itemize}
    \end{column}
  \end{columns}
\end{frame}
